\chapter{Neuronal covariance and the geometry of population responses}\label{ch:population}
\begin{readergoal}Distinguish signal correlation, noise correlation, factor structure, and the information represented in population covariance.\end{readergoal}
\section{One neuron versus a population}
Let $r(t)=(r_1(t),\ldots,r_N(t))^\top$ represent firing-rate estimates for $N$ units. A time series $r(t)$ traces a curve through $\mathbb R^N$. Neurons' responses may change together because of shared inputs, latent states, task drive, or interactions. The covariance matrix captures second-order cofluctuation but does not identify which explanation generated it.

In repeated-stimulus experiments, \emph{signal correlation} measures similarity in mean tuning across stimuli; \emph{noise correlation} measures trial-to-trial cofluctuation after controlling the stimulus or condition. These are distinct statistics, and their estimates depend on temporal window, firing rate, spike sorting, and brain state \cite{CohenKohn2011}.

\section{Shared gain and overdispersion}
A modulated Poisson count model can be written $N\mid G\sim\mathrm{Poisson}(G\mu)$, where $G>0$ is a random gain with $E[G]=1$. Then $E[N]=\mu$ and
\[\Var(N)=\mu+\mu^2\Var(G).\]
The second term represents fluctuations in the latent gain; the decomposition follows directly from the law of total variance. For two neurons with gains $G_1,G_2$, shared gain fluctuations can generate positive count covariance. This does not require a direct synaptic link between their spike generators \cite{Goris2014}.

\section{Latent factors and factor analysis}
A linear observation model $r_t=Cz_t+d+\varepsilon_t$ represents high-dimensional responses through lower-dimensional latent $z_t$. Factor analysis separates low-rank shared covariance from modeled observation noise. GPFA adds smooth temporal covariance priors to latent trajectories. Identifying the effective latent dimensionality requires validation and attention to scaling, stationarity, and neuron subsampling \cite{CunninghamYu2014,Gao2017}.

\section{Random-matrix caution}
For a high-dimensional sample covariance, random eigenvalue fluctuations can be large even when the true covariance is isotropic. In a simple independent-Gaussian null with feature-to-sample ratio $q$, the Mar\v{c}enko--Pastur bulk has support $\sigma^2(1\pm\sqrt q)^2$ under the standard asymptotic assumptions. This is not a universal biological noise threshold; correlated trials, heterogeneous variances, and temporal smoothing violate simple null assumptions.

\begin{exercises}
\item Derive the gain-modulated variance formula using total variance.
\item Explain how simultaneous recording differs from assembling pseudo-populations across sessions.
\item Give a counterexample where noise correlation increases without a change in the direct connection between neurons.
\item Explain why a large first eigenvalue is not automatically a genuine latent neural mode.
\end{exercises}